% 研究ノート: 自己相互作用のある超軽量スカラーの超放射を GWTC-5 の BH スピンで検定する
% ビルド: latexmk -lualatex research_note.tex
% 数値: calc/e0N_*.py をエントリ順に実行(research/ で。共通部品は calc/sr.py ほか)
\documentclass[a4paper,10pt]{ltjsarticle}
\usepackage[margin=22mm,top=22mm,bottom=24mm]{geometry}
\usepackage{amsmath,amssymb,bm}
\usepackage[dvipsnames]{xcolor}
\usepackage{booktabs,tabularx,array,float,xltabular}
\usepackage{enumitem}
\usepackage[most]{tcolorbox}
\usepackage{tikz}
\usetikzlibrary{arrows.meta,positioning,calc,fit,shapes.geometric}
\usepackage{pgfplots}
\pgfplotsset{compat=1.18}
\usepgfplotslibrary{fillbetween,groupplots}
\usepackage[font=small,labelfont=bf]{caption}
\usepackage[colorlinks=true,breaklinks=true,allcolors=MidnightBlue]{hyperref}

% ---------- 記号 ----------
\newcommand{\Msun}{\ensuremath{\,M_\odot}}
\newcommand{\eV}{\ensuremath{\,\mathrm{eV}}}
\newcommand{\GeV}{\ensuremath{\,\mathrm{GeV}}}
\newcommand{\yr}{\ensuremath{\,\mathrm{yr}}}
\newcommand{\Hz}{\ensuremath{\,\mathrm{Hz}}}
\newcommand{\Mpl}{\ensuremath{M_\mathrm{Pl}}}
\newcommand{\arxiv}[1]{\href{https://arxiv.org/abs/#1}{arXiv:#1}}
\newcommand{\review}{LVK 文献レビュー}  % ~/Documents/Claude/LVKtest/Refs/lvk_physics_review/

% ---------- 配色 ----------
\definecolor{cA}{HTML}{2A78D6}
\definecolor{cB}{HTML}{EB6834}
\definecolor{cC}{HTML}{1BAF7A}
\definecolor{cD}{HTML}{8E5BC9}
\definecolor{cM}{HTML}{8A8984}
\definecolor{ink}{HTML}{0B0B0B}
\definecolor{ink2}{HTML}{52514E}
\definecolor{gridc}{HTML}{E4E3DF}
\pgfplotsset{note/.style={
  axis line style={ink2}, tick style={ink2}, tick label style={font=\scriptsize, color=ink2},
  label style={font=\small, color=ink}, major grid style={gridc}, grid=major, clip=true}}

% ---------- 囲み ----------
\newtcolorbox{result}[1][]{colback=cA!5, colframe=cA!60, boxrule=0.5pt, arc=2pt, left=5pt, right=5pt, top=3pt, bottom=3pt, fonttitle=\small\bfseries, coltitle=ink, colbacktitle=cA!14, title={結果}, #1}
\newtcolorbox{caveat}[1][]{colback=cB!5, colframe=cB!70, boxrule=0.5pt, arc=2pt, left=5pt, right=5pt, top=3pt, bottom=3pt, fonttitle=\small\bfseries, coltitle=ink, colbacktitle=cB!15, title={仮定と限界}, #1}

\setlist{itemsep=2pt, topsep=4pt, parsep=0pt}
\numberwithin{equation}{section}

\title{研究ノート\\自己相互作用のある超軽量スカラーの超放射を GWTC-5 の BH スピンで検定する}
\author{}
\date{2026年11月3日開始}

\begin{document}
\maketitle
\vspace{-2.5em}

\begin{tcolorbox}[colback=white, colframe=ink2, boxrule=0.4pt, arc=2pt, left=6pt, right=6pt]
\small
\textbf{このノートの約束}
\begin{itemize}
\item 日付つきのエントリを順に追加する。過去のエントリは書き換えず、訂正は新しいエントリで行う。
\item 数値はすべて \texttt{calc/} のスクリプトから出す(エントリ $n$ は \texttt{calc/e0$n$\_*.py}、出力は \texttt{calc/out/} と \texttt{figs/data/})。
\item 文献の式・数値は arXiv の本文で照合し、式番号を書く。抜き書きは \texttt{Refs/notes/} にある。\review\ にある値は finding の ID(例: T06-06)を示す。
\item 先行研究との比較は、このノートの模型で計算した量で行う。
\end{itemize}
\end{tcolorbox}

\tableofcontents

% エントリ 1(2026-11-03): 解析の設計
\section{[E1, 2026-11-03] 解析の設計}\label{sec:e01}

\subsection{問い}

GWTC-5 の BH スピンの集団は、四次の自己相互作用(強さ $\lambda=\mu^2/f^2$)をもつ超軽量スカラーの質量 $\mu$ と崩壊定数 $f$ を、どこまで排除するか。あるいは支持するか。特に、自由場では両立しない二つの観測事実を、ある範囲の $f$ が両立させるかを決める。
\begin{itemize}
\item \textbf{1 世代目の積み上がり}: 1 世代目の BH のスピンは、質量に依存する臨界値に集まっているように見える。
\item \textbf{2 世代目の速い回転}: 2 世代目らしい BH が $\chi\approx0.7$--$0.8$ で回っている。
\end{itemize}
テーマを選んだ経緯は \texttt{LVKtest/research/topic\_selection\_theory.md}(本文で照合したあとの「訂正」の節を含む)、決めた設計は \texttt{DESIGN.md} にある。

\subsection{現状(本文で照合)}

抜き書きは \texttt{Refs/notes/population\_papers.md} にある。

\paragraph{同じ GWTC-5 から逆の結論}
\begin{itemize}
\item \textbf{排除}(\arxiv{2607.01317}): 257 の合体で $1.7\times10^{-14}\le\mu\le3.3\times10^{-12}\eV$ を排除した(95\%)。
  \begin{itemize}
  \item 自然スピンは質量によらない単一の Beta 分布(同論文の式 2)。
  \item 超放射の率は水素原子近似の解析式で、$m=1$--6 を使う(End Matter の式 4)。
  \item 選択効果は入れていない。
  \item 自己相互作用は QCD アクシオンの $f_a$ でだけ確かめ、最終スピンが「場合により $\sim0.05$ 変わる」とした。ほかの $f_a$ は今後の課題としている(SM の「Superradiance for higher $\alpha$」)。
  \end{itemize}
\item \textbf{証拠}(\arxiv{2609.02678}): $\mu=0.71^{+0.18}_{-0.11}\times10^{-12}\eV$(90\%)、$\ln\mathcal{B}=7.8$(Table I)。
  \begin{itemize}
  \item 模型は、臨界スピン $\chi_\mathrm{SR}(M_f,T_\mathrm{age},\mu)$ を超える確率を、$\chi_\mathrm{SR}$ を中心とする幅 $\sigma_\mathrm{SR}$ の Gauss 分布に寄せる(同論文の式 1)。$\sigma_\mathrm{SR}=0.18^{+0.05}_{-0.03}$。
  \item $\chi_\mathrm{SR}$ は SuperRad で計算する。年齢 $T_\mathrm{age}$ は事象ごとに $[1\,\mathrm{Myr},2.2\,\mathrm{Gyr}]$ で対数一様。
  \item 選択効果は LVK の注入で入れている。
  \item 比べた帰無模型は LVK の Gauss 成分のスピン模型だけで、階層合体の模型はない。自己相互作用は入れていない。
  \end{itemize}
\end{itemize}

\paragraph{分かれる原因}
両者とも自由場である。分かれる原因は積み上がりの幅と、一つの事象である。
\begin{itemize}
\item \arxiv{2609.02678} は自分の模型の幅を 0 にした場合(デルタ関数)も試している。その場合 $3.5\times10^{-14}$--$1.6\times10^{-12}\eV$ がほぼ排除され、それは GW241011 だけによると書いている(End Matter III)。
\item GW241011 の主星は $m_1=19.6^{+3.6}_{-2.5}\Msun$、$\chi_1=0.78\pm0.09$ で、軌道から傾いている。LVK はこれを前の合体の残骸(2 世代目)と解釈している(\review\ T04-28)。
\end{itemize}

\paragraph{空いているところ}
\begin{itemize}
\item LVK のスピンの集団を $(\mu,f)$ で走査した解析はない(2025-01 以降を INSPIRE の被引用検索で確認)。最新のレビュー \arxiv{2610.11091} の \S V も、集団解析は自由場のままだと書いている。
\item 2 世代目の BH の自然スピンを超放射の較正に使った解析もない。
\item 近い仕事:
  \begin{itemize}
  \item \arxiv{2609.12052}: GWTC-5 で $(\mu,1/f)$ 面を扱う。ただし観測量はスピンではなく、銀河核での $q<0.3$ の合体率。
  \item 個々の事象による $(\mu,f)$ の制限: GW231123 と GW190517(\arxiv{2507.20979}、\arxiv{2507.21788})、GW190412 と GW190517(\arxiv{2503.10347})。
  \item 自由場での一つの事象による制限: GW241011(\arxiv{2510.26931})、O4b・O4c の速く回る事象(\arxiv{2608.11620})。
  \end{itemize}
\end{itemize}

\subsection{仮説の形}

自由場では、$\mu\approx7\times10^{-13}\eV$ の雲は、$15\Msun$、$\chi=0.69$ の BH の回転を約 300 年で臨界値まで奪う(E2)。2 世代目の BH が生まれてから次に合体するまでの時間(遅れ)は Myr 以上なので、GW241011 は残らない。

自己相互作用があると、雲は準位間のエネルギー交換と無限遠への放出で準平衡に達し、回転を失う速さは $f^2$ に比例して遅くなる(\arxiv{2011.11646})。すると次のような $f$ の範囲がありうる。
\begin{itemize}
\item 年をとった 1 世代目の BH(遅れは Gyr 程度)は回転を失う。
\item 若い 2 世代目の BH(遅れは 10--100 Myr 程度)は回転を保つ。
\end{itemize}
回転を失う途中の BH が年齢の分布に応じて散らばるなら、\arxiv{2609.02678} が手で入れた幅 $\sigma_\mathrm{SR}$ は物理的に説明できる。この範囲がなければ、その質量はどの $f$ でも排除される。どちらの結論も論文になる。

\subsection{理論の模型}

自己相互作用のある多準位の雲の運動論を解き、最終スピンの写像 $\chi_f(M,\chi_i,\mu,f,t)$ を表にする。
\begin{itemize}
\item \textbf{準位}: まず $211$、$322$、$433$ と、交換で生まれる減衰する準位(\arxiv{2011.11646} の組)。次に $n\le5$ に広げる(\arxiv{2412.03655})。
\item \textbf{過程}: 超放射、準位間の交換(Bose 増強 $\propto N_iN_j(1+N_k)$)、BH への吸収、無限遠へのスカラーの放出、重力波の放出(消滅と準位間の遷移)。
\item \textbf{率}: 超放射の率は、小さい $\alpha$ では解析式(E2)、$\alpha\gtrsim0.2$ では相対論的な率を使う。相対論的な率は SuperRad の当てはめ(\arxiv{2211.03845}、BSD-3)か、\arxiv{2412.03655} の公開コード(Julia、ライセンスの表示なし)と照合する。どちらを使うかはライセンスを確かめてから決める。
\item \textbf{解き方}: 準位間の交換は超放射よりずっと速く、方程式は硬い。占有数の対数を変数にして陰的な解法を使い、準平衡の解析解(\arxiv{2011.11646})と照合する。
\end{itemize}

\subsection{集団の模型}

\begin{itemize}
\item \textbf{二つの世代}: 1 世代目と 2 世代目の混合とする。
  \begin{itemize}
  \item 1 世代目の自然スピンは柔らかい模型にする(Beta 分布と、質量に依存する混合)。
  \item 2 世代目の自然スピンは、数値相対論の残骸の公式から、親の質量比とスピンの分布で決める($\approx0.69$ の周り)。
  \end{itemize}
\item \textbf{年齢}: 1 世代目は星形成から合体までの遅れ、2 世代目は前の合体から次の合体までの遅れとする。それぞれ星団模型の範囲で周辺化する。
\item \textbf{超放射の写像}: 生まれたときのスピンと年齢を、理論の写像で観測時のスピンに移す。
\item \textbf{選択効果と質量}: 選択効果は GWTC-5 の注入で入れる。質量・赤方偏移の模型と同時に推定する。
\item \textbf{比べる模型}: 次の四つを比べる。
  \begin{itemize}
  \item ボソンなし(柔らかい帰無模型。2 世代目を含む)
  \item 自由場
  \item 自己相互作用あり
  \item \arxiv{2609.02678} の模型の再現(このノートの部品で組む)
  \end{itemize}
\end{itemize}

\subsection{データ}

GWTC-5 の事後分布と累積の注入を使う(\texttt{data/README.md})。PBH\_LVKpop と同じファイルなので、置き場所は共用にする方向でユーザーに確かめる。\arxiv{2609.02678} の事後分布(Zenodo 22242056、146 MB)は、照合のために取得を検討する。

\subsection{照合の計画}

\begin{itemize}
\item E2: 自由場の超放射の率を、小さい $\alpha$ の極限と文献の値と照合する。半解析の飽和の扱いを ODE と照合する。
\item E3: 自己相互作用のある運動論。\arxiv{2011.11646} の準平衡の占有数と、回転を失うかどうかの $f$ の境界を再現する。
\item E4: 自由場の集団解析で、\arxiv{2607.01317} の排除と \arxiv{2609.02678} の証拠の両方を、このノートの部品で再現する。違いが積み上がりの幅と GW241011 に帰着することを確かめる。
\end{itemize}

\subsection{完成の条件}

\begin{enumerate}
\item $(\mu,f)$ 面の排除または証拠の地図。QCD アクシオンの線と、ALP として動機のある $f\sim10^{14}$--$10^{17}\GeV$ を描く。
\item \arxiv{2609.02678} の質量が、ある $f$ で生き残るかの判定。生き残るなら、その $f$ の範囲と、予言される重力波の信号。
\item 頑健さ: 遅れの分布、1 世代目の自然スピン模型、2 世代目の割合、波形の系統誤差、GW241011 と GW231123 を入れるか除くか。
\item 自由場の場合に、先行研究の二つの結果をこのノートの部品で再現した比較。
\end{enumerate}

\subsection{日程}

\begin{itemize}
\item 11 月前半: E2(自由場)、E3(自己相互作用の運動論)、写像の表。
\item 11 月後半: データの取得、E4(自由場の集団解析の再現)。
\item 12 月: 自己相互作用ありの集団解析、頑健さ。
\item 2027 年 1 月: 原稿。
\end{itemize}

\begin{caveat}
\begin{itemize}
\item 2 世代目の遅れの下限が、較正の強さを決める。\arxiv{2107.11730} は、2 世代目の多くは若すぎて雲が育たないとした。ただしそれは $\mu=4.1\times10^{-14}\eV$ での話である(同論文 \S4.3)。$\mu\sim10^{-12}\eV$ の $10$--$20\Msun$ では雲は数百年で育つ(E2)。
\item 自己相互作用で回転を失う速さが遅くなる範囲では、雲の放出する重力波とスカラー波も変わる。連続波の探索の制限と矛盾しないかを、結果の $f$ の範囲で確かめる。
\item 先を越される危険がある。2--3 か月で出す。
\end{itemize}
\end{caveat}

% エントリ 2(2026-11-03): 自由場の超放射と 2 世代目の BH
\section{[E2, 2026-11-03] 自由場の超放射: 2 世代目の BH はどこで回転を失うか}\label{sec:e02}

自己相互作用を入れる前の基準として、自由場のスカラーで、生まれたときのスピンが $\chi_0=0.69$ の BH(2 世代目)がいつ、どこまで回転を失うかを計算した。計算は \texttt{calc/sr.py}(共通部品)と \texttt{calc/e02\_freefield.py}、出力は \texttt{calc/out/e02.json}。

\subsection{式}

$\alpha\equiv GM\mu/\hbar c$ とし、$G=c=\hbar=1$ で書く。準位 $|n\ell m\rangle$ の振動数を $\omega_{n\ell m}=E_{n\ell m}+i\Gamma_{n\ell m}$ と書く(\arxiv{1908.10370} の式 2.27)。

\paragraph{エネルギー}
\arxiv{1908.10370} の式 2.28--2.29 を使う。
\begin{equation}
\frac{E_{n\ell m}}{\mu}=1-\frac{\alpha^2}{2n^2}-\frac{\alpha^4}{8n^4}+\frac{f_{n\ell}}{n^3}\alpha^4+\frac{h_\ell}{n^3}\chi m\alpha^5,
\qquad f_{n\ell}=-\frac{6}{2\ell+1}+\frac{2}{n},\quad h_\ell=\frac{16}{2\ell(2\ell+1)(2\ell+2)}.
\end{equation}

\paragraph{成長率}
\arxiv{1908.10370} の式 2.32--2.34 を使う。
\begin{equation}
\Gamma_{n\ell m}=2\tilde r_+C_{n\ell}\,g_{\ell m}\,(m\Omega_H-E_{n\ell m})\,\alpha^{4\ell+5},
\qquad
C_{n\ell}=\frac{2^{4\ell+1}(n+\ell)!}{n^{2\ell+4}(n-\ell-1)!}\Bigl[\frac{\ell!}{(2\ell)!(2\ell+1)!}\Bigr]^2,
\label{eq:e02_rate}
\end{equation}
記号は次のとおり($M=1$ の単位)。
\begin{itemize}
\item $g_{\ell m}=\prod_{k=1}^{\ell}\bigl[k^2(1-\chi^2)+(\chi m-2r_+E)^2\bigr]$。
\item $\tilde r_+=1+\sqrt{1-\chi^2}$、$\Omega_H=\chi/(2r_+)$。
\end{itemize}

\paragraph{規約の注意}
$\Gamma$ は場の振幅の成長率である。占有数 $N$ は $\mathrm{d}N/\mathrm{d}t=2\Gamma N$ で育つ。
\begin{itemize}
\item 小さい $\alpha$ では、式 \eqref{eq:e02_rate} は $\Gamma_{211}\to\mu\alpha^8(\chi-2\alpha\tilde r_+)/48$ になる。占有数の率はその 2 倍で、分母は 24 になる。
\item 本文で照合した \arxiv{1908.10370} は、式 2.32 を $\alpha\lesssim0.2$ で数値解とよく合うとしている。これより大きい $\alpha$ では相対論的な率に替える(E3 以降)。
\end{itemize}

\paragraph{BH の変化}
雲の量子 1 個が BH から $m\hbar$ の角運動量と $\hbar E$ のエネルギーを持ち去る。準位 $|n\ell m\rangle$ は $m\Omega_H(\chi)=E_{n\ell m}$ となる臨界スピンで成長を止める。

\subsection{解き方: 飽和で区切る半解析}

複数の準位を ODE で同時に解くと、飽和したあとの平衡が非常に硬く、発散した(\texttt{DESIGN.md} D3)。そこで次の半解析を使う。
\begin{itemize}
\item 各準位は 1 個の量子から育つ。$e$ 倍の回数 $A_\ell=\int2\Gamma_\ell\,\mathrm{d}t$ が $\ln N_\mathrm{sat}$ に届いたとき飽和する。
\item 飽和すると、BH は臨界スピンまで回転を失う。$N_\mathrm{sat}$ は、臨界スピンまでに渡す角運動量を $m\hbar$ で割った数である。
\item 飽和にかかる最後の数 $e$-fold は無視する($\ln N_\mathrm{sat}\approx180$ に対して数個)。
\item 準位は $211$、$322$、$433$ を入れた。
\end{itemize}
$211$ だけの場合に ODE と比べた(表~\ref{tab:e02_check})。最終スピンの差は最大 $1.4\times10^{-4}$($\alpha=0.15$)である。スピンが中点まで落ちる時刻は、時間の格子の刻み(0.7\%)以内で一致した。

\begin{table}[H]
\centering\small
\caption{半解析と ODE の照合($211$ だけ、$\chi_0=0.69$)。$t_\mathrm{mid}$ はスピンが初期値と最終値の中点まで落ちる時刻。}
\label{tab:e02_check}
\begin{tabular}{@{}rrrcccc@{}}
\toprule
$M\ [M_\odot]$ & $\mu\ [\mathrm{eV}]$ & $\alpha$ & $\chi_f$(ODE) & $\chi_f$(半解析) & $t_\mathrm{mid}$(ODE) & $t_\mathrm{mid}$(半解析) \\
\midrule
15 & $7\times10^{-13}$ & 0.079 & 0.29705 & 0.29705 & 301 yr & 301 yr \\
5 & $1\times10^{-12}$ & 0.037 & 0.14578 & 0.14578 & $4.93\times10^4$ yr & $4.93\times10^4$ yr \\
30 & $3\times10^{-13}$ & 0.067 & 0.25674 & 0.25674 & 2128 yr & 2128 yr \\
60 & $1\times10^{-13}$ & 0.045 & 0.17399 & 0.17399 & $1.28\times10^5$ yr & $1.28\times10^5$ yr \\
10 & $2\times10^{-12}$ & 0.150 & 0.53555 & 0.53569 & 1.83 yr & 1.82 yr \\
\bottomrule
\end{tabular}
\end{table}

\subsection{結果}

\begin{figure}[H]
\centering
% E2: 自由場で 2 世代目の BH(chi0 = 0.69)が回転を失う範囲。データは calc/e02_freefield.py
\pgfplotsset{colormap={spinmap}{color(0)=(cA!85!black) color(0.35)=(cA) color(0.6)=(cA!35) color(0.69)=(white)}}
\begin{tikzpicture}
\begin{groupplot}[group style={group size=2 by 1, horizontal sep=22mm},
  width=0.45\linewidth, height=0.42\linewidth, note]
% (a) t = 1e8 yr での chi_f の地図
\nextgroupplot[xmode=log, ymode=log, xmin=3, xmax=150, ymin=1e-14, ymax=1e-11,
  xlabel={$M\ [M_\odot]$}, ylabel={$\mu\ [\mathrm{eV}]$}, grid=none, axis on top,
  title={\small (a) $t=10^8\yr$ 後の $\chi_f$},
  colorbar, colormap name=spinmap, point meta min=0, point meta max=0.69,
  colorbar style={width=2.5mm, ytick={0,0.2,0.4,0.6}, yticklabel style={font=\scriptsize}, ylabel={$\chi_f$}, ylabel style={font=\scriptsize}}]
  \addplot[matrix plot*, mesh/cols=121, mesh/ordering=x varies, point meta=\thisrow{chif}]
    table[x expr={10^\thisrow{lgM}}, y expr={10^\thisrow{lgmu}}] {figs/data/e02_chif_t8.dat};
  \draw[ink2, line width=0.5pt] (axis cs:10,1e-14) -- (axis cs:10,1e-11);
  \draw[ink2, line width=0.5pt] (axis cs:20,1e-14) -- (axis cs:20,1e-11);
  \draw[ink2, line width=0.5pt, dashed] (axis cs:45,1e-14) -- (axis cs:45,1e-11);
  \addplot[cB, line width=0.8pt, densely dashed, domain=3:150, samples=2] {7.1e-13};
  \node[font=\scriptsize, color=cB, fill=white, fill opacity=0.8, text opacity=1, inner sep=1pt, anchor=south west] at (axis cs:3.3,8e-13) {$7.1\times10^{-13}$\,eV};
% (b) mu = 7.1e-13 eV での最終スピンと GW241011
\nextgroupplot[xmode=log, xmin=3, xmax=150, ymin=0, ymax=0.9,
  xlabel={$M\ [M_\odot]$}, ylabel={$\chi_f$},
  title={\small (b) $\mu=7.1\times10^{-13}$\,eV},
  legend style={font=\scriptsize, at={(0.98,0.02)}, anchor=south east, fill=white, fill opacity=0.85, text opacity=1, draw=none}, legend cell align=left]
  \fill[cM!18] (axis cs:10,0) rectangle (axis cs:20,0.9);
  \fill[cM!18] (axis cs:45,0) rectangle (axis cs:100,0.9);
  \addplot[cA!45, line width=0.9pt] table[x=M, y=chif_t6] {figs/data/e02_regge_mu7.1e-13.dat};
  \addlegendentry{$t=10^6\yr$}
  \addplot[cA, line width=0.9pt] table[x=M, y=chif_t8] {figs/data/e02_regge_mu7.1e-13.dat};
  \addlegendentry{$10^8\yr$}
  \addplot[cA!60!black, line width=0.9pt, dashed] table[x=M, y=chif_t10] {figs/data/e02_regge_mu7.1e-13.dat};
  \addlegendentry{$10^{10}\yr$}
  \addplot[ink2, line width=0.5pt, densely dotted, domain=3:150, samples=2, forget plot] {0.69};
  % GW241011 の主星(LVK 文献レビュー T04-28: m1 = 19.6 +3.6/-2.5, chi1 = 0.78 +/- 0.09、90%)
  \addplot[cB, only marks, mark=*, mark size=1.8pt, error bars/.cd, x dir=both, x explicit, y dir=both, y explicit]
    coordinates {(19.6,0.78) +- (0,0)};
  \addlegendentry{GW241011 主星}
  \draw[cB, line width=0.7pt] (axis cs:17.1,0.78) -- (axis cs:23.2,0.78);
  \draw[cB, line width=0.7pt] (axis cs:19.6,0.69) -- (axis cs:19.6,0.87);
\end{groupplot}
\end{tikzpicture}

\caption{自由場で、$\chi_0=0.69$ の BH が回転を失う範囲。灰色の帯は、GWTC-5 で速く回る部分集団が見つかった質量 $10$--$20\Msun$ と $45\Msun$ 以上(\arxiv{2605.27226})。
(a) $t=10^8\yr$ 後のスピン $\chi_f$(白は回転を失わない)。縦線は質量 10・20・$45\Msun$、橙の破線は \arxiv{2609.02678} の質量。準位 211・322・433 が飽和する範囲が斜めの帯になる($\alpha\propto M\mu$)。
(b) $\mu=7.1\times10^{-13}\eV$ での最終スピン。点は GW241011 の主星(90\%、\review\ T04-28)。}
\label{fig:e02}
\end{figure}

\begin{result}
\begin{itemize}
\item \textbf{$15\Msun$、$\mu=7\times10^{-13}\eV$}: $\alpha=0.079$、占有数の $e$-fold は $1.7\yr$、$\ln N_\mathrm{sat}=179$ である。スピンは約 300 年で中点まで落ち、$0.297$ に止まる。テーマ選定のときの見積もり(e-fold 約 1 年、育ちきるまで $\sim10^2$ 年)は桁で合っていた。
\item \textbf{GW241011 の主星}(\review\ T04-28 の $m_1=19.6\Msun$、$\chi_1=0.78$): $\mu=7.1\times10^{-13}\eV$ の雲があると、スピンは約 35 年で 90\% 区間の下端 0.69 を下回る。$10^{10}\yr$ 後は $322$ も飽和して 0.195 になる。
\item \textbf{帯の中のすべての質量で $\chi_f<0.5$ になる $\mu$}(図~\ref{fig:e02}(a)):
  \begin{center}
  \begin{tabular}{@{}lcc@{}}
  \toprule
  $t$ & $10$--$20\Msun$ & $45$--$100\Msun$ \\
  \midrule
  $10^6\yr$ & $(4.0\text{--}8.9)\times10^{-13}\eV$ & $(1.0\text{--}1.9)\times10^{-13}\eV$ \\
  $10^8\yr$ & $2.4\times10^{-13}$--$1.8\times10^{-12}\eV$ & $(0.60\text{--}3.8)\times10^{-13}\eV$ \\
  $10^{10}\yr$ & $1.4\times10^{-13}$--$2.7\times10^{-12}\eV$ & $(0.35\text{--}5.6)\times10^{-13}\eV$ \\
  \bottomrule
  \end{tabular}
  \end{center}
\item \textbf{2 世代目の較正}: $\mu\approx7\times10^{-13}\eV$ は、$10$--$20\Msun$ の 2 世代目を、遅れが $10^6\yr$ でもすべて回転を失わせる。自由場で \arxiv{2609.02678} の質量を救う道はない。これは同論文が GW241011 について自分で示したこと(E1)と一致する。
\end{itemize}
\end{result}

\begin{caveat}
\begin{itemize}
\item 成長率は解析式で、$\alpha\lesssim0.2$ でしか正確でない。図の上端(大きい $\mu$)の境界は、相対論的な率に替えると動く。
\item 「帯の中のすべての質量」という基準は説明のためのもので、集団解析ではない。
\item 雲の重力波による減衰、伴星による準位の混合、降着は入れていない。どれもスピンを戻さないか(D2)、戻す量が小さい。
\item 2 世代目の自然スピンは 0.69 に固定した。実際は親の質量比とスピンで $0.6$--$0.8$ に広がる。
\end{itemize}
\end{caveat}

% エントリ 3(2026-11-03): 自己相互作用の準平衡による窓の見積もり
\section{[E3, 2026-11-03] 自己相互作用の準平衡: 窓はありうるか(見積もり)}\label{sec:e03}

E1 の仮説を最初に確かめる。問いは、ある $f$ で GW241011 の主星は回転を保ち、1 世代目は回転を失うことがありうるか、である。計算は \texttt{calc/si.py}(共通部品)と \texttt{calc/e03\_quasieq.py}、出力は \texttt{calc/out/e03.json}。式の典拠は \arxiv{2011.11646}(以下 BGLS)で、抜き書きは \texttt{Refs/notes/theory\_rates.md} \S4 にある。

\subsection{式}

占有数を $\varepsilon\equiv N/(GM^2)$ で書き、$\Mpl\equiv1/\sqrt G=1.22\times10^{19}\GeV$、$\lambda=\mu^2/f^2$ とする。二準位($211$、$322$)の主な過程は二つある。
\begin{itemize}
\item $211\times211\to322\times\mathrm{BH}$: $\gamma_\mathrm{BH}=4.3\times10^{-7}\alpha^{11}(\Mpl/f)^4\tilde r_+\mu$(BGLS の式 21、Table V)。
\item $322\times322\to211\times\infty$: $\gamma_\infty=1.1\times10^{-8}\alpha^8(\Mpl/f)^4\mu$(式 18、Table VI)。
\end{itemize}
BGLS の式 34--35 で $\dot\varepsilon_{211}=\dot\varepsilon_{322}=0$ とすると準平衡になる(式 55a,b)。
\begin{equation}
\varepsilon_{211}^\mathrm{eq}=\frac{2}{\sqrt3}\frac{\sqrt{\gamma_\mathrm{SR}\gamma_\infty}}{\gamma_\mathrm{BH}},\qquad
\varepsilon_{322}^\mathrm{eq}=\frac{\gamma_\mathrm{BH}}{2\gamma_\infty}\varepsilon_{211}^\mathrm{eq},\qquad
\dot\chi=-\gamma_\mathrm{SR}\,\varepsilon_{211}^\mathrm{eq}=-\frac{2}{\sqrt3}\frac{\gamma_\mathrm{SR}^{3/2}\gamma_\infty^{1/2}}{\gamma_\mathrm{BH}}\propto f^2.
\label{eq:e03_qe}
\end{equation}
スピンの式は BGLS の式 36 である。$\gamma_\mathrm{SR}$ には、BGLS の主要項の代わりに E2 の率(\arxiv{1908.10370} の式 2.32 の 2 倍)を使う。$322$ の超放射の寄与は $10^{-8}$ 程度なので落とした。

\paragraph{照合}
主要項の率($\kappa_\mathrm{SR}=0.04$)で、BGLS の数値と比べた($\alpha=0.01$、$\chi=0.9$、$f=10^{15}\GeV$)。
\begin{itemize}
\item $\varepsilon_{211}^\mathrm{eq}=0.246$、$\varepsilon_{322}^\mathrm{eq}=6.89\times10^{-6}$ は、式 55a,b の 0.25、$6.9\times10^{-6}$ と合う。
\item 回転を失う時間 $1/|\dot\chi|$ は、$\mu=10^{-12}\eV$ で $2.44\times10^7\yr$ になった。BGLS の式 60 に印刷された $10^7\yr$ は、同論文自身の定数から組んだ値より 2.4 倍小さい。抜き書きの指摘と同じなので、定数から組んだ式を使う。
\end{itemize}

\subsection{比べ方: 年齢の比}

式 \eqref{eq:e03_qe} では回転を失う速さが $f^2$ に比例する。そのため、二つの BH が同じ $f$ で回転を失う時間の比は $f$ によらない。そこで次の二つの時間を比べる($M$ は固定)。
\begin{itemize}
\item $t_\mathrm{2G}$: GW241011 の主星($19.6\Msun$)が $\chi=0.78$ から、90\% 区間の下端 0.69 まで落ちる時間。
\item $t_\mathrm{1G}$: 質量 $M_1$ の 1 世代目が、自然スピン $\chi_0$ から、$211$ の臨界スピン $+\,\delta$ まで落ちる時間。$\delta=0.18$ は \arxiv{2609.02678} の幅 $\sigma_\mathrm{SR}$ にあたる。
\end{itemize}
窓があるには、1 世代目の年齢 $T_1$ と 2 世代目の年齢 $T_2$ が $T_1/T_2\ge R\equiv t_\mathrm{1G}/t_\mathrm{2G}$ を満たさなければならない。

\subsection{結果}

\begin{figure}[H]
\centering
% E3: 窓に要る年齢の比 R = t_1G / t_2G(mu = 0.71e-12 eV)。データは calc/e03_quasieq.py
\begin{tikzpicture}
\begin{axis}[note, width=0.62\linewidth, height=0.42\linewidth, ymode=log, xmin=6, xmax=40, ymin=0.1, ymax=3e3,
  xlabel={1 世代目の質量 $M_1\ [M_\odot]$}, ylabel={$R=t_\mathrm{1G}/t_\mathrm{2G}$},
  legend style={font=\scriptsize, at={(0.98,0.98)}, anchor=north east, fill=white, fill opacity=0.85, text opacity=1, draw=none}, legend cell align=left,
  unbounded coords=jump]
  \fill[cM!18] (axis cs:28,0.1) rectangle (axis cs:40,3e3);
  \node[font=\scriptsize, color=ink2, anchor=north] at (axis cs:34,2e3) {$\alpha>0.15$};
  \addplot[cA, line width=0.9pt] table[x=M1, y=R_c07_d05] {figs/data/e03_ratio_mu0.71.dat};
  \addlegendentry{$\chi_0=0.7$, $\delta=0.05$}
  \addplot[cA, line width=0.9pt, dashed] table[x=M1, y=R_c07_d18] {figs/data/e03_ratio_mu0.71.dat};
  \addlegendentry{$\chi_0=0.7$, $\delta=0.18$}
  \addplot[cC, line width=0.9pt] table[x=M1, y=R_c05_d05] {figs/data/e03_ratio_mu0.71.dat};
  \addlegendentry{$\chi_0=0.5$, $\delta=0.05$}
  \addplot[cC, line width=0.9pt, dashed] table[x=M1, y=R_c05_d18] {figs/data/e03_ratio_mu0.71.dat};
  \addlegendentry{$\chi_0=0.5$, $\delta=0.18$}
  \addplot[ink2, line width=0.5pt, densely dotted, domain=6:40, samples=2, forget plot] {1};
\end{axis}
\end{tikzpicture}

\caption{$\mu=0.71\times10^{-12}\eV$ で、窓に要る年齢の比 $R$。灰色は $\alpha>0.15$ で、二準位の近似が外れる範囲(\arxiv{2412.03655})。}
\label{fig:e03}
\end{figure}

\begin{result}
\begin{itemize}
\item \textbf{GW241011 が許す $f$}($\mu=0.71\times10^{-12}\eV$): 年齢 $T_2$ のあいだ 0.69 を保つ $f$ の上限は、$T_2=10^6,10^7,10^8,10^9\yr$ でそれぞれ $1.4\times10^{13}$、$4.5\times10^{12}$、$1.4\times10^{12}$、$4.5\times10^{11}\GeV$。QCD アクシオンの線($f\approx8\times10^{18}\GeV$)より 6 桁小さい。
\item \textbf{年齢の比}($\mu=0.71\times10^{-12}\eV$、図~\ref{fig:e03}):
  \begin{itemize}
  \item $M_1=15$--$20\Msun$ の 1 世代目は、GW241011 の主星より $R\approx1$--$40$ 倍年をとっていれば回転を失う。窓は閉じていない。
  \item $M_1=8$--$10\Msun$ では $R\approx50$--$900$ が要る。回転を失う速さがおよそ $\alpha^5$ に比例し(式 \eqref{eq:e03_qe} の $\gamma_\mathrm{SR}^{3/2}\gamma_\infty^{1/2}/\gamma_\mathrm{BH}$)、軽い BH ほど遅いからである。
  \item 窓の中では、軽い 1 世代目は回転を保つ。したがって \arxiv{2609.02678} が示した質量とスピンの相関のうち、低質量の側は自由場と違う形になる。
  \end{itemize}
\item \textbf{質量の 90\% 区間の端}: $\mu=0.60$、$0.89\times10^{-12}\eV$ でも同じ傾向になる。$M_1=10\Msun$ での $R$ は、それぞれ $76$--$390$、$18$--$170$ である(\texttt{e03.json})。
\item \textbf{近似の確かめ}(GW241011、$f=4.5\times10^{12}\GeV$):
  \begin{itemize}
  \item $\varepsilon_{211}^\mathrm{eq}$ は自由場で飽和する量の $5.6\times10^{-9}$ 倍で、自己相互作用が成長を止めている。
  \item bosenova の閾値(BGLS の式 88)の 5\% で、届かない。
  \item 準平衡に達するまで約 30 年で、$T_2$ よりずっと短い。
  \end{itemize}
\end{itemize}
\end{result}

\paragraph{意味}
自己相互作用は、自由場では両立しない二つの事実を両立させうる。条件は $f\sim10^{12}$--$10^{13}\GeV$ と、2 世代目が 1 世代目より十分に若いことである。
\begin{itemize}
\item 窓が開くかどうかは、$f$ と、1 世代目・2 世代目の年齢の分布の組で決まる。これは集団解析で決める問いで、E1 の設計がそのまま要る。
\item 年齢の分布の仮定が結論を決めかねないので、遅れの分布を星団模型で周辺化することが論文の中心になる。
\item この窓では、回転を失う途中の BH が年齢に応じて散らばる。\arxiv{2609.02678} が手で入れた幅 $\sigma_\mathrm{SR}$ にあたるものが、物理として出てくる。
\end{itemize}

\begin{caveat}
\begin{itemize}
\item 二準位・主要項の率である。$\gamma_\mathrm{BH}$ は $\alpha\sim0.2$ で「数倍」大きくなる(BGLS 図 21)。相対論的な補正(\arxiv{2412.03655})と $n\le5$ の準位は入れていない。比 $R$ は $\gamma_\mathrm{BH}$ の質量依存性に直接効く。
\item 質量は固定し、BH の質量の変化は無視した(雲は自由場の飽和量の $10^{-8}$ 倍なので小さい)。
\item 1 世代目の判定($\chi_0$、$\delta$)は説明のための選び方で、集団解析ではない。
\item $f\sim10^{12}\GeV$ の ALP が、ほかの観測(連続波・スカラー波の放出、ほかの BH)と矛盾しないかは確かめていない。
\end{itemize}
\end{caveat}


\end{document}
